\documentclass{article}

\usepackage[margin=1in]{geometry}
\usepackage{booktabs}
\usepackage{tabularx}
\usepackage[destlabel]{hyperref}
\usepackage{parskip}
\usepackage{longtable}
\usepackage{array}
\usepackage{enumitem}
\usepackage{pdflscape}

\newcommand{\pass}{{\textbf{Pass}}}
\newcommand{\fail}{{\textbf{Fail}}}

\title{Hazard Analysis\\\progname}

\author{\authname}

\date{}

\input{../Comments.text}
\input{../Common.text}

\begin{document}

\maketitle
\thispagestyle{empty}

~\newpage
\tableofcontents
\newpage
% =====================================================================
\section{Introduction}
% =====================================================================

\subsection{What Is a Cognitive Walkthrough?}

A cognitive walkthrough is a task-based usability inspection method for evaluating how 
learnable an interface is for a first-time user \cite{nng-overview}. Instead of observing 
real users, a cross-functional group of reviewers steps through each action in a predefined 
task and, at every step, answers a fixed set of questions from the point of view of a 
defined user persona. The goal is to find the parts of an interface that are likely 
to derail a new user.

At each step, the faciltator asks the four questions or variations of the four questions. 
Evaluators discuss four questions \cite{nng-overview,nng-workshop}:
\begin{enumerate}
  \item Will the user try to achieve the right result, that is, understand that this step 
  is needed for their larger goal?
  \item Will the user notice that the correct action is available?
  \item Will the user associate the correct action with the result they are trying to 
  achieve?
  \item After the action is performed, will the user see that progress has been made 
  toward the goal?
\end{enumerate}
If the answer to any question is ``no'', the step is recorded as a \fail{}. The group notes 
why, then continues the analysis as though the step had succeeded so that later steps can 
still be examined \cite{nng-workshop}.

The method was first presented by Lewis and colleagues in 1990 for walk-up-and-use interfaces 
such as kiosks \cite{lewis1990}. The streamlined version most teams use today comes from Wharton 
and colleagues \cite{wharton1994}. The process of the cognitive walkthrough relies on reviewer 
expertise instead of participants and works well on early prototypes. The cognitive walkthrough 
is intended to complement usability testing with real users, not replace it \cite{nng-overview}.

\subsection{Purpose of This Report}

This report documents the cognitive walkthrough of the EyeCanDo gameplay wireframes. Following the 
structure recommended by Nielsen Norman Group, it records the methodology, the people and roles involved, 
and the results, with specific problems identified. A summary of results ia also included along with the changes 
planned for the game as a result of the \fail{} recorded in walkthrough notes\cite{nng-overview}.
\pagenumbering{roman}
\end{document}